\documentclass[12pt]{article}
\usepackage[left=2.5cm,top=2.0cm,right=2.5cm,bottom=3.0cm]{geometry}
\usepackage[utf8]{inputenc}
\usepackage[spanish]{babel}
%\usepackage[english, american]{babel}
%\usepackage[spanish,es-tabla]{babel}
\usepackage[linguistics]{forest}
\usepackage{amssymb, amsmath, amsbsy} % simbolitos
\usepackage{longtable} % para tablas largas
\usepackage{graphicx}
\usepackage{fancyhdr}
\usepackage{xcolor}
\usepackage{multirow}
\usepackage{listings}
\usepackage{caption}
\usepackage{subcaption}
%\usepackage{parskip}
\usepackage[skip=12pt plus1pt]{parskip}
\usepackage{pdfpages} % Incluir PDF en documento en LATEX
\usepackage{verbatim} % comentarios
\usepackage{algpseudocode}
\usepackage{algorithm}
\usepackage{pdflscape}
\usepackage{multirow}
\usepackage{afterpage}
\usepackage{array,booktabs,ragged2e}
%\newcolumntype{R}[1]{>{\RaggedLeft\arraybackslash}p{#1}}
\newcolumntype{L}[1]{>{\raggedright\let\newline\\\arraybackslash\hspace{0pt}}m{#1}}
\newcolumntype{C}[1]{>{\centering\let\newline\\\arraybackslash\hspace{0pt}}m{#1}}
\newcolumntype{R}[1]{>{\raggedleft\let\newline\\\arraybackslash\hspace{0pt}}m{#1}}
\floatname{algorithm}{Algoritmo}
\renewcommand{\listalgorithmname}{Lista de algoritmos}
\renewcommand{\algorithmicrequire}{\textbf{Entrada:}}
\renewcommand{\algorithmicensure}{\textbf{Salida:}}
% Comentario con respecto a las referencias:
% Por defecto este documento utiliza el formato IEEEtr para formatear las referencias. Si algun asesor requiere formato APA, solo comente la siguiente linea y descomente la linea debajo. Tambien para que las citas funcionen de manera adecuada, el paquete label debe tener las opciones mostradas, o en su defecto, el año en el listado no lo muestra correctamente \usepackage[english, american]{babel}
% Para utilizar el formato de citas IEEE y comentar los dos parrafos siguientes
\usepackage[backend=bibtex,sorting=none]{biblatex}
% Para utilizar el formato APA, sugiero comentar la linea anterior y descomentar las dos proximas lineas
%\usepackage[backend=biber,style=apa]{biblatex}
%\DeclareLanguageMapping{english}{american-apa}
\makeatletter
\DefineBibliographyExtras{spanish}{%
  \setcounter{smartand}{1}%
  \let\lbx@finalnamedelim=\lbx@es@smartand
  \let\lbx@finallistdelim=\lbx@es@smartand
}
\renewbibmacro*{name:delim:apa:family-given}[1]{%
  \ifnumgreater{\value{listcount}}{\value{liststart}}
    {\ifboolexpr{
       test {\ifnumless{\value{listcount}}{\value{liststop}}}
       or
       test \ifmorenames
     }
       {\printdelim{multinamedelim}}
       {\lbx@finalnamedelim{#1}}}
    {}}
\makeatother
% Estas lineas permiten romper los hipervinculos muy largos en las referencias!!!!
\setcounter{biburllcpenalty}{7000}
\setcounter{biburlucpenalty}{8000}
\addbibresource{x.bib} % ARCHIVO DE BIBLIOGRAFÍA
%\usepackage{url}
\usepackage[bookmarks=true,breaklinks=true,bookmarksopen=false,colorlinks=true,linkcolor=blue]{hyperref}
\usepackage[hyphenbreaks]{breakurl}
% Regla que define explicitamente que caracteres rompen los hipervinculos para separar las lineas
%https://es.overleaf.com/11089898rhgykrqyqytx
% Actualiza en automático la fecha de las citas de internet a la fecha de la compilación del documento
\usepackage{datetime}
\newdateformat{specialdate}{\twodigit{\THEDAY}-\twodigit{\THEMONTH}-\THEYEAR}
%\newdateformat{specialdate}{\twodigit{\THEDAY}-\THEYEAR}
\date{\specialdate\today}
\newcommand{\HRule}{\rule{\linewidth}{0.25mm}}
% CONSTANTES NECESARIAS PARA EL DOCUMENTO ---> MODIFIQUEN A SU CRITERIO
\newcommand{\ncarrera}            {Ingeniería en Tecnologías de la Información}
\newcommand{\nasesorinstitucional}{DR. SAID POLANCO MARTAGÓN}
\newcommand{\NombreAlumno}{ADRIÁN PÉREZ FACUNDO \newline NURIA YARETZI RIVERA ZAMARRÓN}
%Hombres cambien LA por EL
\newcommand{\elolaNombreAlumno}{el}  
\newcommand{\OA}           {o}  %Hombres cambien A por O
\newcommand{\Matricula}           {2130058 \newline 2130010}  %SU MATRICULA
\newcommand{\NombreProyecto}{Guantes traductores de lenguaje de señas}
\newcommand{\fechacarta}{19 de septiembre 2024}
\newcommand{\ncuatrimestre}{Septiembre-diciembre 2024}
\newcommand{\nevalador}{DR. SAID POLANCO MARTAGÓN}
\newcommand{\FechaExposicion}{11 de Agosto de 2021}
\newcommand{\HoraExposionFormatoVenticuatroHoras}{10:00}
\newcommand{\elolaNombreEmpresa}{la}  %Si la empresa es femenino (por ejemplo universidad, usen la) o masculino (el instituto) pongan el
\newcommand{\organismoreceptor}   {Universidad Politécnica de Victoria}
% NOTA: Dos diagonales juntas (\\) indican un saldo de linea. En este caso particular hay 2 (el titulo se ajusta a tres lineas, porque es muy largo. Hacer las adecuaciones pertinentes
\newcommand{\NombreProyectoheader}     {Guantes traductores de lenguaje de señas.}
\newcommand{\nasesorempresaria}   {DR SAID POLANCO MARTAGÓN}
\newcommand{\fechaPortada}               {Septiembre de 2024}
\newcommand{\separacionCorta}{0.0cm}
\newcommand{\separacionLarga}{0.5cm}
\usepackage[overload]{textcase}
\newcommand{\iemph}[1]{\MakeTextUppercase{#1}}
\pagestyle{fancy}
\headheight 45pt
\fancyhead{} % Clear all header fields
\fancyhead[L]{\includegraphics[height=1.00cm]{UTYP.png}}%
\fancyhead[C]{\begin{center}\NombreProyectoheader\end{center}}%
%\fancyhead[R]{\includegraphics[height=1.5cm]{LogoUPV_2019.png}}%
\fancyhead[R]{\includegraphics[height=1.25cm]{LogoUPV_2023.png}}%
\fancyfoot[R]{\thepage} % Clear all footer fields 
\fancyfoot[C]{}
\fancyfoot[L]{}
\DefineBibliographyStrings{english}{%
  references = {Referencias},% replace "references" with "bibliography"  for `book`/`report`
}
\addto\captionsenglish{%
  \renewcommand{\figurename}{Figura}%
  \renewcommand{\tablename}{Tabla}%
} 
\usepackage{wallpaper}
%\renewcommand{\figurename}{Figura}
%\renewcommand{\tablename}{Tabla}
\begin{document}
%-----------------------------------------------------------------------------------------------------------------
% PAGINA 1 - PORTADA
\setcounter{page}{1}
\pagenumbering{roman}
\thispagestyle{empty}
\begin{center}
\begin{tabular}{cp{5cm}c}
\includegraphics[height=2.25cm]{UTYP.png} & 
& \includegraphics[height=2.25cm]{LogoUPV_2023.png}   \\
\end{tabular}
\Large \textbf{UNIVERSIDAD POLITÉCNICA DE VICTORIA}
\vspace{0.5cm}
\hrule
\vspace{0.1cm} 
\hrule
\vspace{0.5cm}
%\HRule \\[\separacionCorta]
\textbf{\iemph{\NombreProyecto}} \\[\separacionLarga]
%\Large \textbf{TESINA}
%\HRule \\[\separacionLarga]
T E S I N A \\
QUE PARA OBTENER EL GRADO DE \\
\textbf{\iemph{\ncarrera}} \\[\separacionLarga]

PRESENTA: \\[\separacionCorta]
%\textbf{\Capitalize{\NombreAlumno}\\[\separacionLarga]
\textbf{\iemph{\NombreAlumno}}\\[\separacionLarga]
%EN CUMPLIMIENTO DE \\[\separacionCorta]
%LA ESTADÍA DE LA CARRERA DE \\[\separacionCorta]
DIRECTOR \\[\separacionCorta]
\textbf{\iemph{\nasesorinstitucional}} \\[\separacionCorta]

CO-DIRECTOR \\[\separacionCorta]
\textbf{\iemph{\nasesorempresaria}} \\[\separacionCorta]

ORGANISMO RECEPTOR \\[\separacionCorta]
\textbf{\iemph{\organismoreceptor}} \\[\separacionLarga]
\end{center}
\begin{flushright}
\iemph{Ciudad Victoria, Tamaulipas, \fechaPortada}
\end{flushright}
\HRule 
% PAGINA 4 - AGRADECIMIENTOS
\clearpage
\ClearWallPaper
\pagestyle{fancy}
%\section*{\centering Agradecimientos}
%\addcontentsline{toc}{section}{Agradecimientos}
%\input{Agradecimientos.tex}
%----------------------------------------------------------------------------------------------------------------
% PAGINA 5 - RESUMEN EN ESPAÑOL
%\clearpage
%\section*{\centering Resumen}
%\addcontentsline{toc}{section}{Resumen}
%\input{Resumen.tex}
%\textbf{Palabras clave:} Tecnología, Sistema web, Desarrollar, Implementar, Empresas, Modulo.
%-----------------------------------------------------------------------------------------------------------------
% PAGINA 6 - RESUMEN EN INGLES
%\clearpage
%\section*{\centering Summary}
%\addcontentsline{toc}{section}{Summary}
%\input{Summary.tex}
%\textbf{Keywords}: Technology, Web system, Develop, Implement, Companies, Module.
%-----------------------------------------------------------------------------------------------------------------
% PAGINA 7 - INDICE
\clearpage
\addcontentsline{toc}{section}{Índice}
\renewcommand\contentsname{Índice}
\tableofcontents
%-----------------------------------------------------------------------------------------------------------------
% CAPITULOS


\clearpage
\pagenumbering{arabic}
\setcounter{page}{1}
%\input{Capitulo1.tex}

\clearpage
%\input{Capitulo2.tex}

\clearpage
\section{Introducción}
La comunicación es una de las actividades más importantes en la vida de cualquier persona, ya que nos permite expresar ideas, sentimientos y conectar con nuestro entorno. Para la mayoría, el lenguaje hablado es el principal medio de interacción, pero para las personas sordas o con dificultades auditivas, la realidad es diferente. En su caso, el lenguaje de señas es la herramienta esencial para poder comunicarse con los demás. Aunque este sistema es completo y eficaz dentro de la comunidad sorda, su principal desafío es que no es ampliamente comprendido por el resto de la población. Esta barrera en la comunicación entre personas sordas y oyentes genera situaciones de incomprensión, que a menudo terminan aislando a quienes dependen del lenguaje de señas. 

El problema de la falta de comunicación efectiva no solo afecta la vida cotidiana, sino también situaciones críticas en sectores como la educación, la atención médica y el ámbito laboral. Si bien existen soluciones tecnológicas como aplicaciones de traducción o servicios de interpretación, no siempre son accesibles o prácticas, ya que muchas dependen de un tercero para su funcionamiento. En este caso se muestra la necesidad de soluciones más innovadoras que faciliten una comunicación fluida sin depender de otros.

Es en este contexto que surge la idea de los guantes traductores de lenguaje de señas, una tecnología diseñada para facilitar la comunicación entre personas sordas y oyentes de manera directa, rápida y eficiente. El proyecto tiene como objetivo desarrollar guantes equipados con sensores que puedan detectar los movimientos de las manos y los dedos, traduciendo estos gestos en tiempo real a texto o voz. Los guantes contarán con sensores flexibles para medir la curvatura de los dedos, así como sensores giroscópicos y acelerómetros que registrarán el movimiento y la orientación de las manos. Todo esto se conectará a través de microcontroladores ESP32, los cuales permitirán procesar y transmitir los datos de manera eficiente.

Una de las principales ventajas de este proyecto es su enfoque en la portabilidad y facilidad de uso. Los guantes estarán diseñados para ser inalámbricos, utilizando conexiones Wi-Fi para transmitir los datos a un dispositivo receptor, donde se llevará a cabo la traducción del lenguaje de señas a voz o texto. Este enfoque aumenta la libertad de movimiento y la comodidad para el usuario, permitiendo que el sistema se utilice en cualquier lugar y momento, sin la necesidad de depender de intérpretes o equipos costosos.

Se planea implementar técnicas avanzadas de aprendizaje automático en el sistema para mejorar tanto la precisión como la adaptabilidad de los guantes traductores de lenguaje de señas. Se utilizará Tiny Machine Learning (TinyML), una rama del aprendizaje automático que se enfoca en modelos ligeros y eficientes, capaces de ejecutarse directamente en dispositivos pequeños y de bajo consumo, como los microcontroladores que controlan los guantes. La incorporación de TinyML permitirá que los guantes no solo realicen la traducción de señas en tiempo real, sino que también aprendan y se ajusten de manera continua a las particularidades del usuario, mejorando su rendimiento a lo largo del tiempo.

El proceso de aprendizaje automático que se integrará en los guantes se basa en el uso de un conjunto de datos específicos de lenguaje de señas. Estos datos consisten en una amplia variedad de gestos y movimientos que las personas sordas utilizan comúnmente para comunicarse. A través de estos datos, el sistema podrá crear modelos de reconocimiento de patrones que detectarán de manera precisa los gestos realizados por el usuario, traduciéndolos a texto o voz.

Uno de los beneficios clave de esta capacidad de aprendizaje es que el sistema podrá adaptarse a las variaciones individuales de cada usuario. Cada persona tiene su propio estilo al realizar señas, que puede variar ligeramente en velocidad, fuerza y ángulo de los movimientos. Esta personalización es crucial, ya que no solo aumenta la precisión de la traducción, sino que también reduce la frustración del usuario al sentir que sus señas están siendo malinterpretadas por el sistema.

El uso de TinyML también contribuye a mejorar la escalabilidad del sistema. Al implementar un enfoque de aprendizaje automático, los guantes podrían ser utilizados en una amplia variedad de contextos sin necesidad de modificaciones significativas en su hardware. Por ejemplo, se podrían personalizar para adaptarse a diferentes lenguas de señas de distintos países, simplemente entrenando el sistema con los gestos y signos correspondientes a cada región. Del mismo modo, podrían integrarse con tecnologías adicionales, como dispositivos de asistencia auditiva, plataformas de comunicación en línea, o aplicaciones móviles, lo que ampliaría aún más su funcionalidad y su impacto en la vida de las personas sordas.

Esta capacidad de aprendizaje continuo no solo hace que el dispositivo sea más útil y adaptable, sino que también lo convierte en una herramienta más inclusiva. El hecho de que el sistema pueda mejorar y adaptarse según las necesidades de cada usuario refuerza la idea de que la tecnología debe ser accesible para todos, independientemente de las diferencias individuales.

Desde un punto de vista técnico, el uso de sensores flexibles y giroscopios permite captar con precisión cada movimiento, lo cual es fundamental para interpretar correctamente el lenguaje de señas, ya que cualquier variación en los gestos puede cambiar el significado. Además, el ESP32 es un microcontrolador potente y económico que permite procesar grandes cantidades de datos en tiempo real, asegurando que la traducción de los gestos sea rápida y eficiente.

A largo plazo, este tipo de soluciones puede generar cambios significativos en la percepción social de la discapacidad auditiva. Los guantes traductores no solo facilitan la comunicación, sino que también envían un mensaje importante sobre la importancia de crear tecnologías accesibles que incluyan a todos. Esto es fundamental en una sociedad que, aunque ha avanzado en términos de inclusión, todavía enfrenta muchos desafíos para garantizar la equidad en la participación de personas con discapacidades. Así, este proyecto no solo tiene un valor técnico, sino también un profundo impacto social, promoviendo un futuro en el que la inclusión y la accesibilidad sean derechos garantizados para todos.

\clearpage
\section{Antecedentes teóricos}
La detección y traducción del lenguaje de señas ha sido un tema de investigación y desarrollo en el ámbito de la tecnología y la inclusión social durante varias décadas. A medida que la tecnología ha avanzado, se han explorado diversas metodologías y herramientas para facilitar la comunicación entre personas sordas y oyentes. Este apartado revisa algunos de los trabajos más relevantes en este campo, centrándose en los avances tecnológicos, las aplicaciones de aprendizaje automático, y las iniciativas que han surgido para abordar el problema de la comunicación en la comunidad sordomuda.

\begin{itemize} 
\item \textbf{Tecnología de Detección de Gestos} \newline
Se centra en interpretar y reconocer los movimientos humanos, siendo fundamental para la traducción del lenguaje de señas. Utiliza diversas metodologías, como la visión por computadora, que analiza imágenes y videos para identificar gestos, y sensores inerciales que capturan el movimiento de las extremidades. Los algoritmos de aprendizaje automático, incluyendo redes neuronales profundas, permiten clasificar gestos a partir de grandes conjuntos de datos, mejorando la precisión del reconocimiento. Sin embargo, enfrenta desafíos como la variabilidad entre usuarios y las condiciones ambientales, lo que exige soluciones robustas para su aplicación en tiempo real. Esta tecnología tiene aplicaciones en interfaces de usuario, asistencia a personas con discapacidades y mejora de la interacción humano-computadora, lo que la convierte en un área de investigación en constante evolución. \newline
Referencia: Koller, O., et al. (2015). "Real-time Gesture Recognition: A Survey". \cite{core_ac_153562887}

\item \textbf{Wearables y Sensores} \newline 
Los dispositivos portátiles han ganado popularidad en el ámbito de la tecnología de detección de gestos, especialmente en el contexto de la traducción del lenguaje de señas. Un ejemplo destacado es el uso de guantes equipados con sensores flexibles que capturan los movimientos de los dedos y las manos. Estos guantes utilizan sensores inerciales y de flexión para detectar la posición y el movimiento de los dedos, permitiendo una traducción más precisa de los gestos en comparación con los sistemas basados en visión. La principal ventaja de los guantes es su capacidad para funcionar en entornos variados sin la necesidad de cámaras, aumentando así su versatilidad y aplicabilidad. Investigaciones como las realizadas por Ranjan et al. (2018) han demostrado que estos guantes pueden detectar gestos en tiempo real, proporcionando una interfaz intuitiva para la comunicación entre personas sordas y oyentes. Este enfoque no solo mejora la calidad de la traducción, sino que también promueve la inclusión social al facilitar la comunicación directa en diversas situaciones de la vida cotidiana.
\newline
Referencia: Ranjan, R., et al. (2018). "A Wearable Glove for Gesture Recognition". \cite{mdpi_sensors_18_367}

\item \textbf{Aprendizaje Automático y Reconocimiento de Patrones} \newline 
El aprendizaje automático ha revolucionado la detección y traducción del lenguaje de señas, ofreciendo soluciones más precisas y adaptables. En el estudio de Wang et al. (2019), se exploran las técnicas de redes neuronales profundas para el reconocimiento de patrones en secuencias de gestos. Estas redes son entrenadas utilizando grandes conjuntos de datos que representan diversos gestos del lenguaje de señas, lo que permite al sistema no solo reconocer gestos individuales, sino también identificar combinaciones y estructuras complejas del lenguaje. Esta metodología ha demostrado mejorar significativamente la precisión de la traducción, lo cual es crucial para una comunicación efectiva entre personas sordas y oyentes. La capacidad de estos modelos para aprender y generalizar a partir de datos variados hace que sean una herramienta fundamental en la búsqueda de soluciones innovadoras para la inclusión social. \newline
Referencia: Wang, Y., et al. (2019). "Deep Learning Techniques for Sign Language Recognition". \cite{mdpi_sensors_18_367}

\item \textbf{Proyectos de Investigación y Aplicaciones Prácticas} \newline 
El proyecto "Sign Language Recognition using Smart Gloves" desarrollado por un equipo de investigadores en la Universidad de Ciudad del Cabo se centra en la creación de guantes inteligentes equipados con sensores para la detección y traducción del lenguaje de señas. Este sistema utiliza sensores de flexión y acelerómetros para captar los movimientos de las manos y los dedos, permitiendo la traducción en tiempo real de los gestos en texto o audio. La interfaz de usuario facilita la comunicación entre personas sordas y oyentes, mejorando la accesibilidad y la inclusión social. Este tipo de proyectos subraya la importancia de la tecnología portátil en la eliminación de barreras comunicativas y en el fomento de una mayor interacción social. \newline 
Referencia: "Sign Language Recognition using Smart Gloves". Universidad de Ciudad del Cabo. \cite{researchgate_sign_language}

\item \textbf{Detección y Traducción de Gestos} \newline 
El artículo "Hand Gesture Recognition Using Machine Learning Techniques: A Review" de Khan et al. (2019) presenta un análisis exhaustivo de diversas técnicas de reconocimiento de gestos de la mano utilizando enfoques de aprendizaje automático. El estudio examina los algoritmos más utilizados, como máquinas de soporte vectorial (SVM), redes neuronales y métodos basados en características, destacando sus ventajas y desventajas en el contexto del reconocimiento de gestos. Además, se discuten los desafíos asociados con la variabilidad de los gestos, el ruido en los datos y la necesidad de conjuntos de datos robustos para el entrenamiento de modelos. Este trabajo subraya la importancia de desarrollar sistemas precisos y eficientes para la traducción de gestos, lo cual es crucial para aplicaciones como la interpretación del lenguaje de señas. \newline 
Referencia: Khan, A. I., et al. (2019). "Hand Gesture Recognition Using Machine Learning Techniques: A Review". \cite{sciencedirect_s187705091931896x}

\item \textbf{Sistemas Basados en Sensores} \newline 
El artículo "Real-time Sign Language Recognition Using Flexible Sensors" de Yang et al. (2020) aborda el desarrollo de un sistema para el reconocimiento de lenguaje de señas en tiempo real utilizando sensores flexibles. Este estudio explora cómo los sensores flexibles pueden capturar los movimientos de las manos y los dedos, proporcionando una solución más precisa y accesible en comparación con los métodos basados en cámaras. Se detalla la integración de estos sensores en dispositivos portátiles, lo que permite una mayor libertad de movimiento para los usuarios. El trabajo destaca la efectividad de este enfoque en entornos del mundo real, así como su potencial para mejorar la comunicación entre personas sordas y oyentes. \newline 
Referencia: Yang, Z., et al. (2020). "Real-time Sign Language Recognition Using Flexible Sensors". \cite{mdpi_sensors_19_2295}


\item \textbf{Aprendizaje Profundo en Lenguaje de Señas} \newline 
El artículo "Deep Learning for Sign Language Recognition: A Review" de Ghosh et al. (2020) presenta una revisión exhaustiva de las técnicas de aprendizaje profundo aplicadas al reconocimiento de lenguaje de señas. Los autores discuten diversas arquitecturas de redes neuronales, como las redes convolucionales (CNN) y las redes recurrentes (RNN), que se utilizan para mejorar la precisión en la detección y traducción de gestos. El estudio resalta cómo el aprendizaje profundo ha revolucionado la forma en que se aborda el reconocimiento de gestos, permitiendo a los modelos aprender de grandes volúmenes de datos y adaptarse a la variabilidad del lenguaje de señas. Además, se analizan los desafíos existentes y se proponen direcciones futuras para la investigación en este campo.  \newline 
Referencia: Ghosh, S., et al. (2020). "Deep Learning for Sign Language Recognition: A Review". \cite{mdpi_applied_sci_10_2458}

\item \textbf{Dispositivos Inteligentes y Tecnología Inalámbrica} \newline 
El artículo "A Smart Glove for Sign Language Recognition: Design and Development" de Sadiq et al. (2021) explora el diseño y desarrollo de un guante inteligente para el reconocimiento de lenguaje de señas. Los autores presentan un sistema que integra tecnología inalámbrica para facilitar la transmisión de datos en tiempo real, lo que mejora la usabilidad y la eficiencia en la comunicación. Se discuten los componentes del guante, incluyendo sensores flexibles y módulos de comunicación, así como el proceso de calibración y entrenamiento del sistema de reconocimiento. Este trabajo destaca la importancia de combinar dispositivos inteligentes con tecnología inalámbrica para crear soluciones accesibles y efectivas en la traducción del lenguaje de señas.  \newline 
Referencia: Sadiq, M. B., et al. (2021). "A Smart Glove for Sign Language Recognition: Design and Development". \cite{sciencedirect_s0957417421003214}

\item \textbf{Proyectos Colaborativos y Comunitarios} \newline 
El artículo "Community-Driven Approaches to Sign Language Recognition: An Insight" de Rodriguez et al. (2021) examina enfoques impulsados por la comunidad para el reconocimiento del lenguaje de señas. Los autores destacan la importancia de involucrar a las comunidades sordas en el proceso de desarrollo de tecnologías de reconocimiento, argumentando que su participación activa no solo mejora la precisión de los sistemas, sino que también garantiza que las soluciones se alineen con las necesidades y deseos de los usuarios finales. Este estudio ilustra varios proyectos donde la colaboración entre investigadores, desarrolladores y comunidades ha llevado a innovaciones significativas en la detección y traducción del lenguaje de señas, promoviendo un sentido de propiedad y empoderamiento en las comunidades afectadas. \newline  
Referencia: Rodriguez, R., et al. (2021). "Community-Driven Approaches to Sign Language Recognition: An Insight". \cite{journaloflanguageandlinguisticstudies}

\end{itemize}


\clearpage
\section{Descripción del problema}

La comunicación es un aspecto fundamental de la interacción humana, y su falta puede resultar en un aislamiento significativo. Las personas sordas enfrentan desafíos únicos en su capacidad para comunicarse con quienes no dominan el lenguaje de señas. A pesar de los avances en la inclusión social, el acceso a la comunicación efectiva sigue siendo un obstáculo para la comunidad sorda. El lenguaje de señas es la principal forma de comunicación para millones de personas en todo el mundo, y su uso no es universalmente reconocido o comprendido por la población oyente. Esta desconexión crea barreras que impiden la interacción social fluida y pueden contribuir a la marginación de las personas sordas en diversas esferas de la vida, incluyendo la educación, el empleo y el acceso a servicios de salud.

La falta de comprensión del lenguaje de señas por parte de las personas oyentes no solo limita la comunicación, sino que también puede generar malentendidos, frustraciones y situaciones de exclusión. Además, en situaciones críticas, como emergencias médicas o desastres naturales, la falta de un medio efectivo para comunicar información vital puede tener consecuencias graves. A menudo, las personas sordas dependen de intérpretes para facilitar la comunicación, pero la disponibilidad de estos profesionales puede ser limitada, lo que agrava la situación. Este hecho resalta la necesidad de soluciones tecnológicas que no solo proporcionen traducción en tiempo real, sino que también sean accesibles y fáciles de usar en la vida diaria.

Los avances en tecnología, particularmente en el ámbito de la inteligencia artificial y los dispositivos portátiles, han abierto nuevas oportunidades para abordar estos problemas. Los guantes inteligentes equipados con sensores flexibles y algoritmos de aprendizaje automático representan una solución prometedora para la traducción del lenguaje de señas. Sin embargo, a pesar de su potencial, la implementación de estas tecnologías enfrenta varios desafíos. La precisión en la detección de gestos, la adaptabilidad a diferentes estilos de comunicación y la aceptación social son solo algunos de los aspectos que deben ser considerados. 

El desarrollo de un sistema que sea práctico y asequible para la comunidad sorda es crucial. La efectividad de cualquier dispositivo dependerá no solo de su funcionalidad técnica, sino también de su aceptación por parte de los usuarios. Esto implica la necesidad de realizar estudios de usabilidad y considerar el feedback de la comunidad sorda durante el proceso de diseño y desarrollo.

La falta de una solución efectiva para la traducción del lenguaje de señas representa un problema significativo que limita la interacción y la inclusión social de las personas sordas. A medida que la tecnología avanza, es esencial explorar y desarrollar dispositivos que no solo superen las barreras de comunicación, sino que también promuevan la equidad y la inclusión en todos los aspectos de la vida. Esta investigación se centra en abordar estos desafíos, con el objetivo de crear un sistema que no solo traduzca el lenguaje de señas de manera efectiva, sino que también fomente una mayor comprensión y empatía entre la comunidad sorda y la población oyente.


\clearpage
\section{Motivación}

La motivación detrás de este proyecto surge de un deseo de contribuir a la sociedad y mejorar la calidad de vida de las personas sordas. En un mundo donde la comunicación es la piedra angular de las interacciones humanas, es inaceptable que una parte significativa de la población se vea limitada por barreras lingüísticas que dificultan su acceso a la información, la educación y la participación plena en la sociedad. A través del desarrollo de guantes inteligentes para la detección y traducción del lenguaje de señas, buscamos no solo ofrecer una solución tecnológica, sino también fomentar una cultura de inclusión y respeto hacia las personas sordas.

El impacto positivo que este proyecto puede tener va más allá de la mera traducción de gestos. Al eliminar la barrera del lenguaje, se facilita la comunicación directa entre personas sordas y oyentes, promoviendo la interacción social y el entendimiento mutuo. Esto no solo mejora la calidad de vida de los usuarios, sino que también enriquece la comunidad en su conjunto al fomentar la diversidad y la empatía. Cada pequeño paso hacia la inclusión contribuye a construir una sociedad más justa y equitativa, donde todos tengan la oportunidad de participar y expresar sus ideas y sentimientos sin restricciones.

Este proyecto tiene el potencial de inspirar a otros a reconocer y abordar las necesidades de la comunidad sorda. Al demostrar que la tecnología puede ser una herramienta poderosa para la inclusión, se abre la puerta a nuevas iniciativas y colaboraciones que pueden transformar la forma en que interactuamos con la diversidad lingüística y cultural. La motivación de este trabajo es, en última instancia, ser un agente de cambio y contribuir a un futuro en el que la comunicación sea accesible para todos, sin importar sus capacidades auditivas. Así, se espera que el desarrollo de estos guantes inteligentes no solo sirva como un avance técnico, sino también como un símbolo de esperanza y progreso hacia una sociedad más inclusiva.


\clearpage
\section{Justificación}

La necesidad de desarrollar guantes inteligentes para la detección y traducción del lenguaje de señas se justifica en múltiples dimensiones. En primer lugar, la comunidad sorda enfrenta desafíos significativos en su vida diaria debido a las barreras de comunicación. A menudo, se ven excluidos de interacciones sociales, oportunidades laborales y acceso a servicios esenciales, lo que resulta en un sentimiento de aislamiento y frustración. La implementación de esta tecnología innovadora busca mitigar esos desafíos, proporcionando a las personas sordas una herramienta que les permita comunicarse de manera más efectiva y fluida con quienes los rodean. Este enfoque no solo responde a una necesidad práctica, sino que también promueve un cambio cultural hacia la aceptación y el respeto de la diversidad.

La justificación técnica del proyecto radica en los avances recientes en tecnologías de sensores, aprendizaje automático y dispositivos portátiles. Estos desarrollos permiten crear soluciones más precisas y accesibles que nunca. Al integrar estas tecnologías en un solo dispositivo, se maximiza la eficiencia del sistema, lo que resulta en una traducción de lenguaje de señas más precisa y en tiempo real. Esto es fundamental, ya que la comunicación efectiva no solo se basa en la traducción literal de gestos, sino también en la interpretación del contexto y la fluidez del lenguaje. 

Por último, la iniciativa tiene el potencial de servir como un modelo para futuras investigaciones y desarrollos en el campo de la inclusión social y la tecnología. La creación de guantes inteligentes no solo beneficiará a la comunidad sorda, sino que también puede inspirar a otros investigadores y desarrolladores a explorar soluciones innovadoras para diversos problemas de comunicación y accesibilidad. Al abordar esta problemática de manera integral y comprometida, se busca no solo ofrecer una solución técnica, sino también contribuir a una sociedad más inclusiva y equitativa donde todas las voces sean escuchadas y valoradas.


\clearpage
\section{Alcances y Limitaciones}

El proyecto de desarrollo de guantes inteligentes para la detección y traducción del lenguaje de señas tiene un alcance definido que incluye varias dimensiones. En primer lugar, se compromete a diseñar, prototipar y probar un dispositivo que permita la traducción en tiempo real de gestos del lenguaje de señas. Este proceso abarcará la integración de sensores flexibles y tecnologías de aprendizaje automático para mejorar la precisión y adaptabilidad del sistema. Además, se prevé la creación de una interfaz de usuario que facilite la interacción con el dispositivo, permitiendo a los usuarios sordos y oyentes comunicarse de manera efectiva.

El proyecto también enfrenta ciertas limitaciones que podrían afectar el cumplimiento de los objetivos planteados. Una de las principales limitaciones es la disponibilidad y calidad del material utilizado en la fabricación de los guantes. Los sensores y componentes electrónicos necesarios para lograr un rendimiento óptimo pueden tener restricciones en cuanto a costo, disponibilidad y calidad. La precisión del sistema dependerá en gran medida de la calidad de estos materiales, lo que podría influir en la eficacia de la traducción del lenguaje de señas.

La complejidad inherente al reconocimiento de gestos es otro factor limitante. El lenguaje de señas es dinámico y varía entre diferentes culturas y contextos, lo que puede dificultar la creación de un modelo de aprendizaje automático que reconozca todos los signos y combinaciones posibles. El entrenamiento del modelo requerirá un conjunto de datos diverso y representativo, lo que puede ser un desafío en términos de recopilación y validación de datos.

El tiempo y los recursos disponibles para la investigación y el desarrollo del proyecto también son factores limitantes. La construcción de un prototipo funcional y su posterior prueba en entornos reales requerirá una planificación cuidadosa y una asignación adecuada de recursos, tanto humanos como materiales. Es importante ser realista acerca de estos desafíos y mantener un enfoque flexible que permita adaptarse a las circunstancias cambiantes a lo largo del desarrollo del proyecto.

%\clearpage
%\addcontentsline{toc}{section}{Índice de figuras}
%\renewcommand\listfigurename{Índice de figuras}

%\listoffigures

%\clearpage
%\addcontentsline{toc}{section}{Índice de cuadros}
%\renewcommand\listtablename{Índice de cuadros}
%\listoftables

%\clearpage
%\addcontentsline{toc}{section}{Índice de algoritmos}
%\renewcommand\listalgorithmname{Índice de algoritmos}
%\listofalgorithms

%-----------------------------------------------------------------------------------------------------------------
% REFERENCIAS
\clearpage
\Urlmuskip=0mu plus 1mu\relax
\addcontentsline{toc}{section}{Referencias} 
\printbibliography
 
\end{document}